\documentclass[11pt]{article}
\usepackage[margin=1in]{geometry}
\usepackage{amsmath,amssymb,amsthm}
\usepackage{booktabs}
\usepackage{hyperref}

\newtheorem{theorem}{Theorem}
\newtheorem{lemma}[theorem]{Lemma}

\title{Exact Coordinates and Proof of Zaman's Recurrence for the Two-Color Spiral Queen Sequence (OEIS A395348)}
\author{Salem Levinstein}
\date{November 26, 2026}

\begin{document}
\maketitle

\begin{abstract}
We prove Rahid Zaman's conjecture on OEIS A395348: in the two-color greedy spiral queen game on $\mathbb{Z}^2$, the Black queen labels satisfy $a(n + 4) = a(n) + 8n + 34$ for every integer $n \ge 2$. Rather than using a large computational cutoff, we establish an exact structural invariant showing that every Black queen has even diagonal coordinates $(x \pm y)$ and $y \in \{0, 1\}$, every Red queen has odd diagonal coordinates and $x \in \{-1, 2\}$, and the projected line sets of both players form exact contiguous intervals at every block $h \ge 1$. This determines exact closed-form coordinates for both A395348 and its companion sequence A395349 for all $n \ge 0$. The entire proof is formally verified in Lean~4.
\end{abstract}

\section{Introduction}

Number the integer lattice $\mathbb{Z}^2$ along the standard counterclockwise square spiral starting at $0 = (0, 0)$ (OEIS A308884 / A274641): ring $r \ge 1$ consists of the $8r$ cells with $\max(|x|, |y|) = r$, starting at label $(2r - 1)^2$ and traversing the East ($e = 0$), North ($e = 1$), West ($e = 2$), and South ($e = 3$) edges of length $2r$ with offset $0 \le t \le 2r - 1$:
\begin{align*}
  e = 0 &: (r,\, 1 - r + t), &
  e = 1 &: (r - 1 - t,\, r), \\
  e = 2 &: (-r,\, r - 1 - t), &
  e = 3 &: (1 - r + t,\, -r),
\end{align*}
at spiral label $k = (2r - 1)^2 + 2re + t$.

Two players, Black and Red, alternate turns in rounds $n = 0, 1, 2, \dots$. On round $n$, Black places a queen at the smallest unoccupied label $a(n)$ not attacked by any existing Red queen (sharing none of $x$, $y$, $x + y$, $x - y$), and Red then places a queen at the smallest unoccupied label $b(n)$ not attacked by any existing Black queen. The Black sequence $a(n)$ is OEIS A395348~\cite{OEIS-A395348}, and the Red sequence $b(n)$ is OEIS A395349. On November 18, 2026, Rahid Zaman conjectured that $a(n + 4) = a(n) + 8n + 34$ for all $n \ge 2$.

\begin{theorem}\label{thm:main}
Let $B(n)$ and $R(n)$ denote the $n$-th Black and Red placements. For $0 \le n \le 3$,
\[
  B(0) = (0,0),\quad B(1) = (1,1),\quad B(2) = (-2,0),\quad B(3) = (3,1),
\]
\[
  R(0) = (2,-1),\quad R(1) = (-1,2),\quad R(2) = (-1,-2),\quad R(3) = (2,3).
\]
For every $h \ge 0$ and $j \in \{0, 1, 2, 3\}$, setting $n = 4h + 4 + j$:
\begin{center}
\begin{tabular}{@{}ccccc@{}}
\toprule
$j$ & Ring $r$ & Black $B(n)$ & Red $R(n)$ & Black label $a(n)$ \\
\midrule
$0$ & $2h + 3$ & $(-r, 1)$ & $(2, -r)$ & $4r^2 + r - 1$ \\
$1$ & $2h + 4$ & $(r, 0)$ & $(-1, r)$ & $4r^2 - 3r$ \\
$2$ & $2h + 4$ & $(-r, 0)$ & $(-1, -r)$ & $4r^2 + r$ \\
$3$ & $2h + 5$ & $(r, 1)$ & $(2, r)$ & $4r^2 - 3r + 1$ \\
\bottomrule
\end{tabular}
\end{center}
Consequently, $a(n + 4) = a(n) + 8n + 34$ for every integer $n \ge 2$.
\end{theorem}

\section{Proof of Theorem~\ref{thm:main}}

Because occupied and opponent-attacked cell sets only grow, if a player places a queen at label $k^*$, no label $\le k^*$ can be legal on any subsequent turn of the same player. Hence on each turn it suffices to verify that every cell strictly between the player's previous queen and the predicted target is attacked by an existing opposing queen, and that the predicted target is neither occupied nor attacked.

Direct evaluation of rounds $0 \le n \le 7$ (all within ring $r \le 5$, label $\le 93$) yields
\begin{align*}
  a(0..7) &= (0, 2, 18, 28, 38, 52, 68, 86), &
  b(0..7) &= (9, 15, 21, 31, 47, 61, 75, 93),
\end{align*}
matching Theorem~\ref{thm:main}. For $h \ge 1$, let $\mathcal{B}_h = \{B(i) : 0 \le i \le 4h + 3\}$ and $\mathcal{R}_h = \{R(i) : 0 \le i \le 4h + 3\}$, and write $X(S)$, $Y(S)$, $D_+(S)$, $D_-(S)$ for the projections of $S$ onto $x$, $y$, $x + y$, $x - y$.

\begin{lemma}\label{lem:inv}
For every $h \ge 1$:
\begin{align*}
  X(\mathcal{B}_h) &= [-(2h+2),\, 2h+3] \setminus \{-1, 2\}, & Y(\mathcal{B}_h) &= \{0, 1\}, \\
  D_+(\mathcal{B}_h) &= 2[-(h+1),\, h+2], & D_-(\mathcal{B}_h) &= 2[-(h+1),\, h+1], \\
  X(\mathcal{R}_h) &= \{-1, 2\}, & Y(\mathcal{R}_h) &= [-(2h+2),\, 2h+3] \setminus \{0, 1\}, \\
  D_+(\mathcal{R}_h) &= 2[-(h+2),\, h+2] + 1, & D_-(\mathcal{R}_h) &= 2[-(h+2),\, h+1] + 1.
\end{align*}
\end{lemma}

\begin{proof}
At $h = 1$, direct enumeration of rounds $0 \le n \le 7$ verifies all eight identities. Passing from $h$ to $h + 1$ adjoins the four Black points $(-(2h+3),1)$, $(2h+4,0)$, $(-(2h+4),0)$, $(2h+5,1)$ and the four Red points $(2,-(2h+3))$, $(-1,2h+4)$, $(-1,-(2h+4))$, $(2,2h+5)$, extending each interval by the exact required endpoints while preserving $Y(\mathcal{B}_h) = \{0, 1\}$ and $X(\mathcal{R}_h) = \{-1, 2\}$.
\end{proof}

In particular, every Black queen has $x \pm y$ even and $y \in \{0, 1\}$, while every Red queen has $x \pm y$ odd, $x \in \{-1, 2\}$, and $y \notin \{0, 1\}$. Thus no predicted target is ever occupied by an earlier queen (same-color labels strictly increase, and the Black and Red $y$-sets are disjoint), and no target is ever attacked diagonally by the opposing color.

\begin{proof}[Proof of Theorem~\ref{thm:main}]
Assume Theorem~\ref{thm:main} holds through round $4h + 3$ ($h \ge 1$). Using Lemma~\ref{lem:inv} and updating the four line sets after each placement, we check all eight turns of block $h$:
\begin{enumerate}
  \item \emph{Black, $j = 0$ ($r = 2h + 3$):} Gap from $(r, 1)$ to $(-r, 1)$ has $2 \le y \le r = 2h + 3$, all in $Y(\mathcal{R}_h)$. Target $(-r, 1)$ has $x \notin \{-1, 2\}$, $y = 1 \notin Y(\mathcal{R}_h)$, and $x \pm y$ even.
  \item \emph{Red, $j = 0$ ($r = 2h + 3$):} Black now has $X = [-r, r] \setminus \{-1, 2\}$. Every gap cell from $(2, r)$ to $(2, -r)$ has $x \in X$ except $(-1, r)$ (killed by $x + y = r - 1 \in D_+$) and $(-1, -r)$ (killed by $x - y = r - 1 \in D_-$). Target $(2, -r)$ has $x = 2 \notin X$, $y = -r \notin \{0, 1\}$, and $x \pm y$ odd.
  \item \emph{Black, $j = 1$ ($r = 2h + 4$):} Red now has $Y = [-(r - 1), r - 1] \setminus \{0, 1\}$. Every gap cell from $(-(r - 1), 1)$ to $(r, 0)$ has $y \in [-(r - 1), -1] \subset Y$ except $(-(r - 1), 0)$ (killed by $x + y = -(r - 1) \in D_+$). Target $(r, 0)$ has $x \notin \{-1, 2\}$, $y = 0 \notin Y$, and $x \pm y$ even.
  \item \emph{Red, $j = 1$ ($r = 2h + 4$):} Black now has $X = [-(r - 1), r] \setminus \{-1, 2\}$. Every gap cell from $(2, -(r - 1))$ to $(-1, r)$ has $x \in X$ except $(2, r)$ (killed by $x - y = -(r - 2) \in D_-$). Target $(-1, r)$ has $x = -1 \notin X$, $y = r \notin \{0, 1\}$, and $x \pm y$ odd.
  \item \emph{Black, $j = 2$ ($r = 2h + 4$):} Red now has $Y = [-(r - 1), r] \setminus \{0, 1\}$. Every gap cell from $(r, 0)$ to $(-r, 0)$ has $y \in [2, r] \subset Y$ except $(r, 1)$ ($x + y = r + 1 \in D_+$) and $(-r, 1)$ ($x + y = -(r - 1) \in D_+$). Target $(-r, 0)$ is legal.
  \item \emph{Red, $j = 2$ ($r = 2h + 4$):} Black now has $X = [-r, r] \setminus \{-1, 2\}$. Every gap cell from $(-1, r)$ to $(-1, -r)$ has $x \in [-r, -2] \subset X$. Target $(-1, -r)$ is legal.
  \item \emph{Black, $j = 3$ ($r = 2h + 5$):} Red now has $Y = [-(r - 1), r - 1] \setminus \{0, 1\}$. Every gap cell from $(-(r - 1), 0)$ to $(r, 1)$ has $y \in [-(r - 1), -1] \subset Y$ except $(r, 0)$ ($x + y = r \in D_+$). Target $(r, 1)$ is legal.
  \item \emph{Red, $j = 3$ ($r = 2h + 5$):} Black now has $X = [-(r - 1), r] \setminus \{-1, 2\}$. Every gap cell from $(-1, -(r - 1))$ to $(2, r)$ has $x \in X$ except $(2, -(r - 1))$ ($x + y = -(r - 3) \in D_+$). Target $(2, r)$ is legal.
\end{enumerate}

Finally, substituting Black's coordinates for $n = 4h + 4 + j$ into $k = (2r - 1)^2 + 2re + t$ yields the four label formulas in Theorem~\ref{thm:main}. Advancing $n \mapsto n + 4$ increases $r$ by $2$, giving $a(n + 4) - a(n) = 16r + 18 = 8n + 34$ for $j \in \{0, 2\}$ and $a(n + 4) - a(n) = 16r + 10 = 8n + 34$ for $j \in \{1, 3\}$. Together with $a(6) - a(2) = 68 - 18 = 50 = 8 \cdot 2 + 34$ and $a(7) - a(3) = 86 - 28 = 58 = 8 \cdot 3 + 34$, this establishes $a(n + 4) = a(n) + 8n + 34$ for all $n \ge 2$.
\end{proof}

\section{Formal Verification}

Theorem~\ref{thm:main} (exact closed-form coordinates for both colors, uniqueness of the greedy game, and Zaman's recurrence for all $n \ge 2$) has been formally verified in Lean~4 (module \texttt{ZamanA395348}) directly from the greedy spiral rules, using only the standard axioms (no \texttt{sorry}, no custom axioms, and no \texttt{native\_decide}). An independent Python verifier (\texttt{verify.py}) simulates 408 rounds (816 placements) and confirms all invariants and the recurrence.

\begin{thebibliography}{9}
\bibitem{OEIS-A395348}
OEIS Foundation Inc., \emph{The On-Line Encyclopedia of Integer Sequences}, Sequence A395348 (conjecture by R.~Zaman, November~18, 2026), \url{https://oeis.org/A395348}.
\end{thebibliography}

\end{document}
